\documentclass{report}
\usepackage{amsmath,amssymb}
\setlength{\parindent}{0mm}
\setlength{\parskip}{1em}
\begin{document}
\begin{center}
\rule{6in}{1pt} \
{\large Ralph Thesen \\
{\bf Efficent adaptive solution of multilevel-systems}}

Fraunhofer Institute for Algorithms and Scientific Computing SCAI \\ Schloss Birlinghoven \\ D-53754 Sankt Augustin \\ Germany
\\
{\tt ralph.thesen@scai.fraunhofer.de}\end{center}

The Gauss-Southwell-algorithm is the greedy grandfather of the
well known Gauss-Seidel-algorithm. This adaptive algorithm solves
in every step the equation with the largest current residual
instead of a priori fixed running order. We present a method to
find the maximum residual with minimum costs. We employ a heap
data structure to store the residual which allows us not only to
find the maximum residual very fast, but also to update entries
in an incremental manner.

The greedy choice of the largest residual leads to a fast
convergence, since the iterative solver adapts itself to the
structure of the linear system. The beneficial effects from
utilizing the structural information in a greedy manner dominate
the extra costs in a total balance. In particular multilevel
systems, which are highly structured, are well suited for such an
adaptive algorithm.

This adaptive strategy shows a dramatic speedup for anisotropic
problems compared with the non-adaptive methods. Convergence
rates are shown for a class of highly anisotropic linear elliptic
problems. In all examples our proposed method is substantially
faster than the Gauss-Seidel-algorithm.


\end{document}
